\documentclass[12pt,english]{reviewresponse}

%% Language
\usepackage{babel}
\usepackage[babel]{microtype}
\usepackage[babel]{csquotes}
\usepackage{xcolor}

%% Fonts
\usepackage[T1]{fontenc}
\usepackage{lmodern}
%\usepackage{newcent} % different font
%\usepackage[scaled]{beramono} % different monospace font


%% Bibliography
\usepackage[backend=biber,style=ieee,dashed=false,url=false,isbn=false,defernumbers=true,refsection=section]{biblatex}
\bibliography{literature.bib}

\usepackage{hyperref}

\newcommand{\DC}[1]{{\color{violet}[{\tiny D:~}#1]}}


\title{Quantifying Very Extreme Precipitation and Temperature Using Huge Ensembles Generated by Machine Learning-based Climate Model Emulators}
\author{Christopher Paciorek and Daniel Cooley}
\journal{BAMS}

\begin{document}
\maketitle


% Response to Editor
\editor
\begin{generalcomment}
	Reviewers 1 and 2 raise important concerns about the overall focus of the paper, and Reviewers 2 and 3 question the use of the emulator. It is essential that these overarching issues are addressed in your revision.
\end{generalcomment}
\begin{revresponse}

There are two primary concerns here: first, whether EVA is actually needed given the ability to generate very large ensembles (particularly in the future), and second, the importance of our focus on how to do statistical analysis compared to whether the emulator is fit for purpose.

Regarding the first point, we agree that it could be 
the case at some point in the future that one will be able to use empirical quantiles to "directly" estimate AEP values (i.e., return levels) even for very-low probability events. When that will occur is unknown, and there are good reasons to think that it may be some time in the future:
\begin{itemize}
    \item This work uses an emulator operating at low spatial resolution (1 degree). For many applications, higher resolution will be desired and will presumably be much more computationally expensive.
\item Similarly, extremes for shorter time durations (e.g., 1 hour) will be desired for some applications (or for analyzing 24-hour precipitation rather than daily precipitation). This work uses an emulator running at 6-hour time steps, and higher resolution will presumably be much more computationally expensive.
\item For climate change scenarios, one will need to address nonstationarity, either in a continuous fashion with the statistics of extremes modeled statistically as changing over time or via running multiple stationary scenarios under different conditions. This will require larger emulator sample sizes. 
\item As discussed in relation to reviewer 1 comments, if one wanted to do seasonal analysis, this would involve smaller sample sizes and might require larger ensembles.
\end{itemize}
Furthermore, high-end GPUs needed for emulation are expensive (e.g., NVIDIA's H200 GPU costs roughly \$30,000 currently). Despite the fact that ML ensembles can be generated with much less computational time, there are and will continue to be financial barriers prohibiting researchers from generating huge ensembles.  
The ability to estimate very high quantiles empirically is still far off on the horizon.

Our purpose in studying {\em this} ensemble is that it provides a testbed to assess EV methods for quantifying extremes that few, if any, other current products allow.
As mentioned by reviewer 1, stochastic weather generators (SWGs) could be used to produce an equivalent number of years of simulated weather.
However, a SWG is produced by fitting a statistical model to weather, and then drawing from this statistical model, and the tail behavior of the produced weather is determined by the model fit to the data.
In contrast, the emulator (being trained on ERA5 output, predicting one time step ahead) does not have known tail behavior, so investigating how well EV methods extrapolate into the tail is a much more authentic investigation than if applied to SWG output.
There have been efforts to produce huge ensembles from output from physically-based models; e.g., LENS2\footnote{https://www.cesm.ucar.edu/community-projects/lens2}.
However, as LENS2 is a transient run, one would have to consider nonstationarity before using its output for a similar investigation.
Multi-model experiments can also be used to produce large ensembles, but one would need to account for model-to-model variability in order to use a multi-model ensemble to assess EV extrapolation.
As recognized by the reviewers, this ensemble's space/time resolution is inadequate for estimation of many risk-relevant quantities, and further, it is untested whether this ensemble is able to produce extreme events with realistic intensity and frequency.
Nevertheless, this ensemble's size (amount of output) does provide an authentic testbed to assess EV methods that we argue are an important part of quantifying extremes that can have major impacts on society and ecosystems. We have added text at the end of the introduction (lines 122-127) to justify our choice to use the ML-based emulator.


We also want to point out that empirical quantiles are also statistical estimators, with associated uncertainty. E.g., with 101 observations, one estimates the 99th percentile as the 2nd-largest observation. If one took a different sample of data, that 2nd-largest observation would be different. Suppose that we are interested in the 1-in-100000 year AEP value for precipitation and instead of our 10000-year ensemble, we were able to produce a one-million year ensemble using new computational resources at some point in the future. Our estimate would be, roughly, the 10th-largest observation, and there would be non-negligible uncertainty. If one were interested in the 1-in-one-million year AEP value, the estimate would be, roughly, the 2nd largest observation, with even more uncertainty. Statistical analysis would still be needed to quantify uncertainty.
Compared to the empirical estimator which uses only this $k$th-largest-order statistic, EVA uses information from all observations above a (suitably high) threshold, and thus EVA may produce estimates with less uncertainty than using empirical quantiles. 
This work does not investigate the question of using empirical quantiles versus EVA, as we only have a 10000-year ensemble, but the point is that even much larger ensembles still raise statistical questions. Furthermore, the statistical uncertainty as quantified by a standard error decreases as the square root of the ensemble size, so improvements in certainty require ever larger sample sizes.


As a final motivation, NOAA has been charged with producing updated PMP estimates in the coming years, so there is a need for estimating AEP values now and not just many years in the future when EVA might no longer be needed. 

We have added text to the discussion (lines 400-404) stating why we think that improved computational efficiency is unlikely to obviate the need for EVA in the near future.

Regarding the second primary concern, we agree that fitness-for-purpose is a critical issue, and we think we are clear in pointing this out while also saying that that is a topic that we expect will be explored in depth by other studies. Within the length constraints of BAMS, we can't do a careful investigation of fitness and indeed that is not the focus of our work. As presented in our BAMS paper proposal, this work focuses on the statistical aspects of using massive emulator output, leaving emulator assessment to the various ongoing efforts in the community. Given the accumulating evidence regarding the usefulness of ML-based emulators (and the massive investments driving related AI/ML technology forward), we expect that at least some emulators will be found to reasonably emulate the real world (or emulate climate models) and therefore be useful. An important question then is how best to analyze the emulator output to quantify extremes. That is the focus of our work. 


\end{revresponse}

\reviewer 

\begin{revcomment}
    The authors present an extreme value analysis (EVA) of a large ensemble climate model emulator simulation of precipitation and surface temperature. The focus is squarely on the value of peaks-over-threshold (POT) EVA methods for estimating very rare quantiles. The paper is well-written and I applaud the authors for emphasizing the most challenging aspect of the problem, which is not the statistics but rather the dynamics—i.e., how does one generate and validate rare events in a weather/climate model, whether that model is an emulator or a fully numerical model. That said, I do have some comments and concerns about the paper, and I must admit I am torn over whether BAMS is the right destination for this paper, given the relatively limited scope of the work. I defer to the other reviewers and ultimately the editors to decide that.
\end{revcomment}

\begin{revresponse}
Please see our general response to the editor.
\end{revresponse}
\begin{revcomment}
    

The usage of the emulator in this work feels like a bit of a red herring, despite the authors being very clear about its limitations. There are other ways of generating thousands of years of plausible precipitation and temperature series, including at much higher resolution than the emulator. The most obvious would be a stochastic weather generator, many of which (specifically ones that don't jointly model space and time dimensions) can be deployed quickly and easily in a large number of locations and can provide watershed or point (e.g. station-scale) estimates with seasonally varying parameters, including at subdaily scales. These results, by virtue of their better resolutions, would likely constitute a more meaningful starting point for an EVA than these very spatially coarse ACE2 emulator results. Of particular concern here is that the coarse scale of the emulator results means that analysis of tail heaviness can be quite divorced from real world situations at many
decision-relevant scales, and that lessons learned from lighter tails (as the authors note in their work) don't transfer well to fine scales. I can appreciate that the authors perhaps wanted to use an emulator since such tools seem to be en vogue these days, but they don't make a fully convincing case that the emulator is bringing very much to the table, given its multiple limitations and given that they don't focus much on understanding its capabilities.
\end{revcomment}

\begin{revresponse}

We agree with the reviewer that the coarse spatial and temporal resolution of this emulator means that it cannot be used to produce estimates of very low probability events at decision-relevant space-time scales.  We believe ML emulators will increase their resolution in years to come and will produce output at scales that are relevant for decision-making.  Whether a future ML emulator's output accurately characterizes very low probability, very extreme events is an essential question that will require extensive investigation by interdisciplinary teams of statisticians and atmospheric scientists.

We disagree, however, that if performed on SWG output rather than this emulator output, our EVA investigation would yield the same insights.  As explained in our response to the editor above, since SWG output is generated from a fitted statistical model, it has the tail behavior prescribed by that fitted model.  Given complete knowledge of the fitted model, any high quantile could be solved for analytically with no estimation required.  In contrast, the tail behavior of the emulator is unknown, and investigating how well an EVA's estimates of the emulator's tail compare to empirically-derived quantile estimates obtainable from the massive ensemble is authentic in a manner that applying to SWG output would not be. 



\end{revresponse}



\begin{revcomment}
On a related note, if an emulator is cheap to run, do we need EVA at all? I don't know what 36 hours on an A100 costs, but presumably in a context in which an emulator can provide continental-scale results for moderate computational effort, it would be feasible to run the emulator for 360 hours or 3600 hours, given the task can be broken into pieces. With 100,000 years or more of simulations, maybe we don't need to fit theoretical distributions, POT or otherwise. That would have clear advantages, because a quantile estimate from an EV distribution doesn't provide any of the critical "event-scale" information (i.e., spatiotemporal rainfall structure) that is crucial in flood estimation and that would in principle be available from empirical quantile estimates (since an empirical quantile can be associated directly to a specific storm event).
\end{revcomment}

\begin{revresponse}
Please see our general response to the editor.  

We disagree with the statement that an EV analysis does not provide "event-scale" information.
With emulator output, one has spatio-temporally consistent fields (to the extent the emulator captures the real-world dependence structure). One can then aggregate a variable (e.g., precipitation) to any space-time scale of interest. At that point, one could apply EVA or one could use empirical quantiles. And one can look at the space-time structure of individual storms associated with the largest precipitation events in the emulator output regardless of whether one is using empirical quantiles or EV methods. So we don't see that the ability to quantify "event-scale" information differs depending on whether one is using EVA or quantiles. 

And building on our response to the previous comment, stochastic weather generators would not generally provide high-quality space-time fields (at least without heroic statistical assumptions about space-time dependence structure). In contrast, with an emulator, the dependence emerges as a result of emulating the evolution of the system over time.  



\end{revresponse}

\begin{revcomment}
The authors could provide a potential  counterargument by emphasizing that data-driven and particularly dynamical methods that produce results at more decision-relevant resolutions (such as the "kilometer scale" identified in the NASEM report) would likely be vastly more computationally expensive than existing emulators and thus the EVA element may prove necessary for computational expediency.
\end{revcomment}
\begin{revresponse}
Yes, we agree and see our general response to the editor and added text in the resubmission.
\end{revresponse}

\begin{revcomment}
  L81-82: Since BAMS is read by a general audience, more care should be taken to define what is meant by "point estimate", since this means something quite different for meteorologists and hydrologists than for statisticians
  
\end{revcomment}

\begin{revresponse}
We think the reviewer is using "point estimate" to mean an estimate at a particular point in space (although in our Google searching there seems to be quite a bit of use in meteorology/hydrology of "point estimate" in the statistical sense of a single estimate without uncertainty). In any event we decided to simply use "statistical estimate" (line 79) and avoid the issue of what 'point' means.


\end{revresponse}
\begin{revcomment}
    
L215: This was a bit confusing but basically it seems like the authors are using 83 years of (annual?) maxima from ERA5 to compare against ACE2 ensemble predictions. This seems a bit problematic, since while it's likely ok to assume stationarity within the 22-year emulator training period, stationarity almost certainly doesn't exist within the length of the full ERA5 record. Does that influence the interpretation of this result? I am struggling to figure that out myself, since as my next comment says, I don't quite understand what Figure 1 shows.
\end{revcomment}

\begin{revresponse}
The main point of Fig.\ 1 is to show that the extremes become increasingly larger as the ensemble size increases, with the ERA5 values serving as an approximate frame of reference. This allows us to answer the first-order question of "can an emulator produce extremes larger than those in the data on which it was trained?"  

Yes, the figure could be interpreted to assume stationarity. The emulator values could in principle be higher than the ERA5 values on average because the emulator was run for 2001-2022, while the ERA5 data include much earlier years. However Fig.\ 1a actually provides no evidence of this, since the emulator and ERA5 values are generally similar (along the 1:1 line). 

In any event, we replotted Fig.\ 1a and Fig.\ 7a using ERA5 data only for 2001-2022 (to match the emulator forcing data) and the results were qualitatively very similar, although of course the spread of points is larger because the ERA5 values are more variable when based on fewer years) (see Response Fig.\ 1 above).

\begin{figure}
    \centering
    \includegraphics[width=\linewidth]{prec-temp-compare-era5-2001-2022.pdf}
    \caption{Comparison of maximum (a) precipitation [compare to Manuscript Fig.\ 1a] and (b) temperature [compare to Manuscript Fig.\ 7a] by grid point for ERA5 (2001-2022) and 22 years of emulator output forced with 2001-2022 values.}
\end{figure}

 

\end{revresponse}

\begin{revcomment}
 Figure 1: I'm a bit confused. Do you mean annual max? Or the actual maximum over the whole time period? And each data point represents a grid cell.
   
\end{revcomment}
\begin{revresponse}
Each point in the plot represents the maximum daily value over 83 years (for ERA5) or 100/1000/10000 years (for the ensemble) for a given grid cell. We have tried to clarify this by revising the Figs. 1 and 7 captions to say "Maximum daily precipitation for each grid cell over the specified time period". 

\end{revresponse}
\begin{revcomment}
Figure 2: Could that high value in Kansas in panel e be a numerical stability issue?
    
\end{revcomment}
\begin{revresponse}
Good question -- we checked the fit for the Kansas location. The optimization converged successfully (including getting the same result with different starting values) and nothing stands out apart from the shape parameter estimate being quite high (as seen in Fig.\ 2c), which leads directly to high estimated uncertainty for very low AEP depths. There are other fairly high shape parameter estimates in a swath of grid points from western Texas through Wisconsin, so this one location is only somewhat anomalous. We also note that in EVA fitting, one tends to see that the shape and scale parameters trade off (e.g., estimates are negatively correlated), with a ridge in the likelihood surface. Indeed, the Kansas location has a shape estimate larger than its neighbors and a scale estimate smaller than its neighbors (not shown). The AEP depth estimate itself is not extreme relative to its spatial neighbors and the "outlying" standard error is not large compared to the depth estimate. The viridis color scheme here is standard but perhaps a bit misleading in terms of the yellow looking more extreme than it is. So this appears to simply be an effect of sampling variability. If desired in real analysis, situations like this or more problematic such sitations could be addressed using various forms of statistical regularization.

% NOTE: for our reference of the top 20 values, only one pair occur on adjacent (or even nearby) days.
\end{revresponse}

\begin{revcomment}
L249-250: The block maxima estimation exploration would be useful very pretty easy to do! If it were me, I'd include it.
    
\end{revcomment}

\begin{revresponse}
We didn't focus on block maxima for increasing block sizes in part for space reasons and in part because it felt more natural to us to increase the threshold than to increase the block size as a way to focus the analysis on the most extreme observations, since even with larger blocks, some less extreme values could occur by chance. Under statistical asymptotics the two are equivalent, as we note in the text. We also argue that it may be easier for a climate scientist to pick a threshold based on scientific understanding and the mixture of storm types at a location than to pick a block size. E.g., with the latter one would need to think about things like "how often would an extreme hurricane occur here" rather than simply looking at the distribution of precipitation values. 


We generated a plot analogous to Fig.\ 3 but for the GEV with increasing block sizes (see Response Fig.\ 2 below). The results are very similar to Fig.\ 3. Given this and space limitations, we have simply commented on this result (lines 251-253) and not provided an additional figure in the manuscript.

\begin{figure}
    \centering
    \includegraphics[width=\linewidth]{prec-rv-bias-gev-only.pdf}
    \caption{AEP depth estimates for each grid cell from GEV fit with increasing block lengths (columns, left to right) compared to AEP depth estimates from empirical quantiles.}
\end{figure}

\end{revresponse}


\begin{revcomment}
    
L281-282: I'm not really convinced by this argument, that moving away from storm type/seasonality is a feature and not a bug. To me this implies a move away from "understanding," which seems unhelpful. From a practical standpoint, estimation of probable maximum flood does require understanding of the seasonal/weather/"on-the-ground" context in which the extreme precipitation takes place. I'm not saying that POT methods should be avoided, but rather that this aspect—being able to ignore sources of complexity—is not, in my mind, where the attraction/benefit lies.
\end{revcomment}


\begin{revresponse}
We are not arguing for moving away from understanding; rather we expect that quantitative AEP estimates of the form we investigate would be used in conjunction with physical understanding and context. E.g., one might estimate extreme AEP depths and also examine the simulated fields in the emulator for the locations and times where the extreme precipitation occurred to better inform use of the AEP depth estimates. 

That said, in retrospect, our wording of "ignore" is indeed rather strong and perhaps misleading, so we see why the reviewer reacted negatively to our framing. In lines 286-288 we have reworded our characterization, noting that avoiding the need for storm typing when estimating quantiles may be a helpful simplification, given that assigning types can be a difficult exercise and one with associated uncertainty. 

We also point to text in lines 385-398 (also present in the original submission) that highlights the usefulness of examining the extreme events produced by the emulator.

Also note that if one did season-specific or storm-type specific estimation, sample sizes will be smaller, so there is all the more need for EVA.

\end{revresponse}



\begin{revcomment}
L294: I may be missing something, but I don't see any explanation in the Methods section of how standard errors or other statistical uncertainty characterization is carried out.
    
\end{revcomment}

\begin{revresponse}
The use of the delta method to estimate uncertainty is mentioned in lines 152-153 (formerly lines 151-152 in the original submission). We've added "on the MLE-based information matrix" as well.
\end{revresponse}

\begin{revcomment}
  L315-318 and Fig 8: that local maximum in southern Minnesota doesn't look physical to me, and invoking sampling variability here isn't convincing. Can you trace it back to the emulator, or might this be the MLE optimization freaking out?
  
\end{revcomment}
\begin{revresponse}

The ``hotspot" in southern Minnesota does not appear to arise from obvious outliers in the emulator temperature values (the threshold exceedances). The MLE optimization converges. The MLE estimate for the shape parameter is unusually high (Fig.\ 8c), resulting in a somewhat high estimate for the AEP temperature estimate (Fig.\ 8b) and the high (but not very large in absolute terms -- standard error of about 3 and estimate of about 50) uncertainty (Fig.\ 8e).
As described in our response to Comment 8, estimates like this arise due to the negative correlation in the shape and scale parameters of EV distributions.
The spatial extent of the estimates (best shown in Fig.\ 8c) is due to the spatial correlation of daily temperature data (which has a longer-range correlation than precipitation data).
We agree that the large positive shape parameter is probably not physical.
Again, a more comprehensive analysis could employ a form of regularization to reduce the variability of the shape parameter estimates. 

We have added a bit more text on this to our original text and moved the text to the Fig.\ 8 caption.

% Note: of top 20 temp values in the grid cell at middle of hotspot, 4 occur in one heatwave, 3 in another, 2 in a third, and 3 in a fourth, so temporal autocorr is a factor here; say that actual analysis would remove repeated values from the same event; however when I did this for that grid cell, it didn't have much impact.

\end{revresponse}


\reviewer 
\begin{revcomment}
   Reviewer \#2: The authors state in the abstract that "a recent National Academies report on modernizing PMP estimation proposed using very large climate model-based ensembles to estimate extreme quantiles, possibly through machine learning-based ensemble boosting," and that "here we assess such an approach…". However, later in the abstract they also write, "our focus is on how to use huge ensembles to estimate very extreme statistics…". These statements appear somewhat contradictory and suggest the authors are uncertain about what story to tell in their paper.
 
\end{revcomment}

\begin{revresponse}
Our story is specifically about the statistical aspects of analyzing large ensembles, likely generated by ML-based emulators (as we do here), but not necessarily.  

Our wording of "here we assess such an approach" is indeed not as explicit as it could be. We have changed the wording in line 16 to be "here we assess the *statistical* aspects of the approach suggested in NASEM report in the context of emulator-based ensembles". 
\end{revresponse}

\begin{revcomment}
    Given the large sample size (10,560) and the capability of machine learning tools to generate arbitrarily large ensembles, I question the necessity of focusing on "how to use huge ensembles to estimate very extreme statistics." With a machine learning-based approach, it should be feasible to generate a 100,000-year equivalent sample. With such a large sample, one should be able to robustly estimate a one-in-10,000-year event using a straightforward empirical quantile method. Thus, the core issue is not the statistical inference of rare events per se, but whether a machine learning model trained on a limited dataset can produce a large ensemble that is valid for estimating such rare events.

\end{revcomment}

\begin{revresponse}
    Please see our general response to the editor, where we give an explanation on why we believe that using empirical quantile estimates from ML ensembles will not be feasible in the near future and why EV methods will continue to be necessary.  Additionally, 
    note our discussion of the fact that empirical quantiles are estimators with uncertainty and that with a 100000-year ensemble the AEP value for 1-in-10000 years is, roughly, the tenth-largest value, which will likely have non-negligible sampling variability, because it is so far in the tail. 
\end{revresponse}


\begin{revcomment}
    
In this context, I find limited value in publishing a paper like this in a high-impact journal such as BAMS. Focusing on "how to use huge ensembles…" seems like a missed opportunity. There are large ensembles produced by multiple GCMs. Some of the large ensembles are also dynamically downscaled by RCMs. These large ensembles already provide testbed for evaluating the approach proposed in the National Academies report. For example, one could use a single ensemble member to train the machine learning model and then generate a large synthetic ensemble. This process could be repeated across multiple ensemble members. If the machine learning-generated rare event statistics are consistent across these synthetic ensembles and align with those from the original large ensemble, it would build confidence in the approach. Conversely, if the results vary significantly, it suggests that the machine learning model cannot reliably extrapolate beyond the information contained in the
limited training set. Regardless of whether the outcome is positive or negative, such an investigation would be a meaningful contribution to the field.
\end{revcomment}
\begin{revresponse}
Please see our general response to the editor.

More specifically, we could instead have carried out our analysis on large ensembles from GCMs or from RCM-based downscaling, although there are stationarity issues and model-to-model differences that could complicate using physically-based model output as the testbed (see our general response for more discussion of this). Our primary focus is on the statistical aspects, and the emulator provides a convenient tool to create an ensemble. Emulators are seeing widespread use and good performance, including for climate as well as weather, so we think assessing the statistical methods in that context is well-justified.

The reviewer seems to be suggesting an approach to compare results across different GCMs to assess consistency. That's outside the scope of our work, which is to consider how to analyze large ensembles statistically. 
\end{revresponse}



\reviewer

Reviewer \#3: (Review by John Nielsen-Gammon) This paper analyzes a massive AI-based ensemble to explore the use of extreme value statistical analysis in quantifying events corresponding to very large return periods, a technique that has been previously recommended but not tested.   I have a few minor comments and other flagged wording issues.  Comment numbers refer to line numbers in the text. 

Minor comments:

\begin{revcomment}
 50: Insert "location-specific" after "empirical"?  Distinguishes approaches such as radar-based transposition, which holds promise in estimating magnitudes of extremely rare events.
   
\end{revcomment}

\begin{revresponse}
We've inserted "location-specific" (line 50).
   
\end{revresponse}

\begin{revcomment}
   79-80: You can say one thing relevant about a 1-in-10000 year event: its return value is bounded from below by the return value of the lowest-probability event for which you CAN make a meaningful estimate.
 
\end{revcomment}

\begin{revresponse}
This is of course true, although it's not clear how useful that bound would be, particularly in a context like PMP.
\end{revresponse}

\begin{revcomment}
  106-108: I don't see the distinction.  A block maxima approach can also exclude certain seasons, and blocks can be set to be long enough so as to exclude (almost all) events that statistically do not produce the desired return interval event magnitude.  Also implying block maxima can't do this are lines 121-122.
  
\end{revcomment}

\begin{revresponse}
Yes, setting long blocks is an alternative to using a very high threshold. We say this directly in line 104 (formerly lines 105-106). As we argue elsewhere, using a threshold directly excludes smaller values, while using longer blocks does not, although they are asymptotically equivalent.

As far as lines 121-122, one might read that as implying block maxima can't do this, but read in context after the previous paragraph, we feel this is clear. We slightly reworded the bullet point to read "Can an appropriate threshold-based analysis avoid ...". Also, as noted in response to reviewer 1 and your later comment, we have clarified our presentation of the POT and GEV results to make clear their equivalence and that one could use block maxima provided the block lengths are sufficiently large (lines 249-253).
\end{revresponse}

\begin{revcomment}
 150: I don't know what is meant here by "extrapolation".  Even the smallest sample should include ten or so events that exceed 1-in-1000 year AEP values, so estimating those values empirically requires interpolation rather than extrapolation.
   
\end{revcomment}


\begin{revresponse}
Yes, we agree this is not very clear. Rather than spending precious words on trying to clarify, we have simply removed the sentence. 
\end{revresponse}

\begin{revcomment}
  Fig.\ 2ff: In this and subsequent figures, indicate the units in the figure or caption.
  
\end{revcomment}
\begin{revresponse}
    We have added units to the figures where they were missing (Figs. 2abde, 8abde). 
\end{revresponse}

\begin{revcomment}
	237-250: The GEV comparison should be deleted or made more rigorous.  As it is, the deck is stacked against GEV analysis because of the much larger number of samples in the single GEV analysis compared to even the largest sample size in the POT analysis.  The text strongly implies that the difference in behavior is inherent in POT vs. GEV, but without an apples-to-apples sample size comparison it is impossible to attribute this to the type of analysis rather than the sample size.  The second and third columns of Fig.\ 3 suggest that a POT analysis with a sample size of 10,560 values might look similar to the GEV analysis with 10,560 values.  There should be at least one comparison with similar sample sizes, or else don't include a GEV analysis at all.
\end{revcomment}

\begin{revresponse}
The reason for our focus on annual maxima was that blocks of one year are by far the most common approach to GEV analysis and would likely be the default for one analyzing emulator data. Our submission makes the important point that this is a bad idea for precipitation. In lines 105-106 of the original submission we said that longer blocks are an alternative. In line 244, we said that our conclusion about GEV analysis is under the condition of "based on annual maxima", and in lines 248-249 we stated that one could instead use longer block.

That being said, we have now also carried out GEV analysis with longer block sizes (also suggested by reviewer 1 and see Response Fig.\ 2), which confirms that one could instead use the GEV with longer blocks as an alternative to using a high threshold. As noted in response to reviewer 1, we comment on this in the revision (lines 251-253), but we did not include an additional figure given space and financial limitations. 

As you hypothesize, a POT analysis with sample size equivalent to annual-maxima-based GEV gives very similar results between POT and GEV (see Response Fig.\ 3 above for a version of Fig.\ 3 from the submission that uses POT with sample size of n=10560 instead of annual-maxima-based GEV). For space reasons we have not added an additional figure, but we note this result in lines 250-251.

 \begin{figure}
    \centering
    \includegraphics[width=\linewidth]{prec-rv-bias-pot-only.pdf}
    \caption{Alternative version of Manuscript Fig.\ 3 using POT with n=10560 instead of annual maxima-based GEV.}
\end{figure}

\end{revresponse}


\begin{revcomment}
    

265-269: The main concern with using an AI emulator for this purpose is that the emulator may not be able to produce the proper number of extreme events that exceed the events in its training database.  The paper tests this visually and qualitatively in Fig.\ 1, indicating that the AI emulator can produce at least some events that exceed the training events.  But there's still the concern that it may not produce enough such events.  How would that manifest itself quantitatively?  Presumably, as one increases the POT threshold closer to or beyond the realm of historical events, the AI database would exhibit an unrealistic truncation in intensity, which would manifest itself as a change in the shape parameter from that of an unbounded distribution toward that of a bounded distribution. This is precisely what Figure 4 (bottom row) shows happening, yet this connection is not made in the text.  The text simply notes that the emulated data could be unrealistic, without noting that
this evidence is directionally specific toward a negative answer to the first key question of the paper: is the emulator fit for purpose?
\end{revcomment}

\begin{revresponse}
The fact that the shape parameter estimates appear to have a downward trend as the threshold increases is a very interesting finding.  The reviewer is right to question whether this apparent trend is realistic; that is, whether we would find such a trend in observations (were they available to assess) or in output from a physical model.  This negative trend could in fact arise from an inability of (current?) machine-learning-based emulators to generate data that correctly extrapolates extremes beyond what the emulator is trained on.  
We do note that in going from Fig.\ 1a to 1b to 1c, there does not seem to be some cutoff that precipitation values do not exceed and values in 10560 years of emulator output are much higher than seen in ERA5.
Still, the reviewer's comment centers around a key question that would need to be answered for emulator output to be used for the purpose of producing estimates of very extreme events:  How can we verify that the emulator {\em correctly} extrapolates beyond the data on which it is trained?  This is a key area for joint research between the statistics and atmospheric science research communities.  Fig.\ 1 answers only the question ``can the emulator extrapolate?" not whether it does so correctly. 

That being said, we also note that the apparent need for a high threshold (or long blocks) is consistent with EV theory and with the Ben-Alaya et al. paper we reference. And we think that, properly caveated, it is a finding worth presenting in the literature that future work can investigate further.

\end{revresponse}


\begin{revcomment}

274-277: Determining the overall AEP depth as the largest among the four seasons produces a low-biased estimate of the overall AEP depth.  While this bias would be small in many or most cases, why would you do it this way at all, rather than producing an estimate derived from the combined seasonal probabilities?  For example, if the extreme events come equally from spring and summer, the AEP depth estimated from your technique for an 0.001 probability event would in reality correspond to an 0.002 probability event.
\end{revcomment}

\begin{revresponse}
This is a good point and an oversight on our part. We have modified how we do this to carry out the simple optimization needed to find the AEP value at which the sum of the probabilities over the four seasons equals the desired probability, updated the plots, and modified our wording of how we do this (lines 279-280). In any event, as the reviewer suspected, the bias is small for precipitation and essentially non-existent (as expected) for temperature, so the overall qualitative results (and the appearance of the figures) do not change. 
\end{revresponse}

\begin{revcomment}
 286: Clarify here that you are specifically focusing on statistical (aleatoric) uncertainty.  Also, in the absence of a known criterion for maximum uncertainty for practical use, the discussion should be framed as "how much" the uncertainty can be reduced rather than "whether" it can be reduced enough.
\end{revcomment}

\begin{revresponse}
We have added "statistical" to clarify that is our focus (line 292). 

As far as "how much" versus "whether", while we agree on the absence of known criteria, we do think that users may be able to decide in some situations whether the statistical uncertainty is sufficiently small and that in other situations, the sample sizes may be large enough that statistical uncertainty would be judged to be negligible. Given that, we have chosen to keep the wording as is.
\end{revresponse}

\begin{revcomment}
   Consider estimating the return values for one or both of two famous recent very extreme events: Hurricane Harvey (area-averaged rainfall) and the Pacific Northwest heat wave of 2021.  These can be compared to published estimates.
 
\end{revcomment}
\begin{revresponse}
    This is a good idea from a general scientific perspective and something we had considered including. However, given space limitations and our desire to focus on the statistical questions, we have chosen not to include this. One concern of ours is that this would be seen as considering fitness-for-purpose and would raise questions about why we did not do other analyses related to that question.
\end{revresponse}


\begin{revcomment}
 22: It's not clear which part of the sentence the word "necessary" is referring to.
   
\end{revcomment}
\begin{revresponse}
We've modified the placement of "necessary for precipitation" to clarify this (lines 21-22).
   \end{revresponse} 
\begin{revcomment}
    39: The analysis doesn't "provide" an approach, and anyhow the point of the paper is not to present an approach (we already know about extreme value analysis) but to demonstrate that it is potentially viable.
\end{revcomment}


\begin{revresponse}
We have changed "provides" to "supports" (line 39).

\end{revresponse}

\begin{revcomment}
52, 53, 56: Describes PMP as a "concept" and as an "approach", then says PMP has been estimated using an "approach".  I think it's a concept that uses an approach, but it's not an approach itself.
\end{revcomment}
\begin{revresponse}
We've changed "approach" to "concept" to be consistent (line 52).
\end{revresponse}

\begin{revcomment}

56-58: The description leaves off the critically important step known as moisture maximization.
\end{revcomment}
\begin{revresponse}
We've reworded the sentence to include transposition, maximization and orographic adjustment, which is how the NASEM report summarized PMP as well (line 58).     
\end{revresponse}
\begin{revcomment}

67: The NASEM study only recommended the use of climate models if and when they are found to be fit for purpose, so mention that.
\end{revcomment}
\begin{revresponse}
   Yes, it would be reasonable to include that caveat, but we haven't found a way to insert that without interrupting the flow of the sentence, and since essentially any use of climate models comes with the need to consider fitness-for-purpose, we don't see a specific need to insert that consideration here, particularly as we highlight the question of fitness-for-purpose elsewhere.
 
\end{revresponse}



\begin{revcomment}

74: Typically much less than 5%.
\end{revcomment}
\begin{revresponse}
    We have inserted "much" (line 72).
\end{revresponse}

\begin{revcomment}

423: renormalized block maxima
\end{revcomment}

\begin{revresponse}
 Thanks, good catch. We have inserted "block" (line 433).
\end{revresponse}

\end{document}